% Part: first-order-logic
% Chapter: proof-systems
% Section: seqeunt-calculus

\documentclass[../../../include/open-logic-section]{subfiles}

\begin{document}

\iftag{FOL}
      {\olfileid[fr]{fol}{prf}{seq}}
      {\olfileid[fr]{pl}{prf}{seq}}

\olsection{Le calcul des séquents}

Alors que de nombreux systèmes de !!{derivation} manipulent des
agencements d'!!{sentence}s, le calcul des séquents manipule des
\emph{séquents}. Un séquent est une expression de la forme
\[
!A_1, \dots, !A_m \Sequent !B_1, \dots, !B_n,
\]
c'est-à-dire un couple de suites d'!!{sentence}s, séparées par le
symbole de séquent~$\Sequent$. Chacune des deux suites peut être vide.
Une !!{derivation} dans le calcul des séquents est un arbre de séquents
dont ceux qui figurent tout en haut ont une forme particulière (on les
appelle « séquents initiaux » ou « axiomes »), tandis que tout autre
séquent se déduit de ceux qui se trouvent immédiatement au-dessus de
lui par l'une des règles d'inférence. Ces règles manipulent les
!!{sentence}s des séquents (en les ajoutant, en les supprimant ou en
les réordonnant, à gauche ou à droite), ou introduisent une !!{formula}
complexe dans la conclusion de la règle. Par exemple, la règle
$\LeftR{\land}$ permet de passer de $!A, \Gamma \Sequent \Delta$ à
$!A \land !B, \Gamma \Sequent \Delta$, et la règle $\RightR{\lif}$
de $!A, \Gamma \Sequent \Delta, !B$ à
$\Gamma \Sequent \Delta, !A \lif !B$, quels que soient $\Gamma$,
$\Delta$, $!A$ et~$!B$. En particulier, $\Gamma$ et~$\Delta$ peuvent
être vides.
\begin{editorial}
Note de l'édition française : dans l'expression générale du séquent,
l'indice final de la suite de droite est noté $n$ au lieu de $m$ dans
l'original. Les deux suites ont des longueurs indépendantes.
\end{editorial}

La relation $\Proves$ fondée sur le calcul des séquents se définit
ainsi : $\Gamma \Proves !A$ si et seulement s'il existe une suite
$\Gamma_0$ dont chaque terme $!A$ appartient à~$\Gamma$, et une
!!{derivation} ayant le séquent~$\Gamma_0 \Sequent !A$ à sa racine.
$!A$ est un théorème du calcul des séquents si le séquent
$\Sequent !A$ admet une !!{derivation}. Par exemple, la
!!{derivation} suivante établit que $\Proves (!A \land !B) \lif !A$ :
\begin{prooftree}
  \Axiom$!A \fCenter !A$
  \RightLabel{\LeftR{\land}}
  \UnaryInf$!A \land !B \fCenter !A$
  \RightLabel{\RightR{\lif}}
  \UnaryInf$\fCenter (!A \land !B) \lif !A$
\end{prooftree}

Un ensemble $\Gamma$ est incohérent dans le calcul des séquents s'il
existe une !!{derivation} de $\Gamma_0 \Sequent$, où chaque
$!A \in \Gamma_0$ appartient à~$\Gamma$ et où le membre droit du
séquent est vide. La règle \RightR{\Weakening} permet de
!!{derive} n'importe quel !!{sentence} d'un ensemble incohérent.

Le calcul des séquents a été inventé dans les années 1930 par Gerhard
Gentzen. Son organisation systématique et symétrique en fait un
formalisme très utile pour élaborer une théorie des !!{derivation}s.
Il est relativement facile d'y trouver des !!{derivation}s, mais celles-ci
sont souvent difficiles à lire, et leur rapport avec les démonstrations
n'est pas toujours facile à discerner. Le calcul des séquents s'est
néanmoins révélé une approche très élégante des systèmes de
!!{derivation}, et de nombreuses logiques possèdent un tel calcul.

\end{document}
